\documentclass[14pt,a4paper]{extarticle}

\usepackage{fontspec}
% Проверка наличия шрифтов и установка fallback-вариантов
\IfFontExistsTF{Times New Roman}{\setmainfont{Times New Roman}}{\setmainfont{DejaVu Serif}[Scale=0.92]}
\IfFontExistsTF{DejaVu Sans Mono}{\setmonofont{DejaVu Sans Mono}}{}

\usepackage{polyglossia}
\setdefaultlanguage{russian}
\setotherlanguage{english}
\usepackage{geometry}
\geometry{left=30mm,right=10mm,top=20mm,bottom=20mm}
\usepackage{setspace}
\onehalfspacing
\usepackage{indentfirst}
\setlength{\parindent}{1.25cm}
\setlength{\parskip}{0pt}

\usepackage{graphicx}
\usepackage{float}
\usepackage{caption}
\captionsetup[figure]{justification=centering,labelsep=endash,name=Рисунок}
\captionsetup[table]{justification=raggedleft,singlelinecheck=false,labelsep=endash,name=Таблица}

\usepackage{titlesec}
\usepackage{tocloft}
\usepackage{fancyhdr}
\usepackage{lastpage}
\usepackage{hyperref}
\usepackage{fvextra}
\usepackage{enumitem}
\usepackage{microtype}

\hypersetup{hidelinks}
\sloppy

\pagestyle{fancy}
\fancyhf{}
\fancyfoot[C]{\thepage}
\renewcommand{\headrulewidth}{0pt}
\renewcommand{\footrulewidth}{0pt}

\titleformat{\section}{\normalfont\bfseries\centering}{\thesection}{1em}{}
\titleformat{\subsection}{\normalfont\bfseries\centering}{\thesubsection}{1em}{}
\titleformat{\subsubsection}{\normalfont\bfseries\centering}{\thesubsubsection}{1em}{}
\titlespacing*{\section}{0pt}{18pt}{12pt}
\titlespacing*{\subsection}{0pt}{14pt}{10pt}
\titlespacing*{\subsubsection}{0pt}{12pt}{8pt}

\renewcommand{\cfttoctitlefont}{\hfill\normalfont\bfseries}
\renewcommand{\cftaftertoctitle}{\hfill}
\renewcommand{\contentsname}{СОДЕРЖАНИЕ}
\renewcommand{\refname}{СПИСОК ИСПОЛЬЗУЕМОЙ ЛИТЕРАТУРЫ}
\setcounter{tocdepth}{3}
\setcounter{secnumdepth}{3}

\DefineVerbatimEnvironment{Code}{Verbatim}{breaklines=true,breakanywhere=true,fontsize=\scriptsize,frame=single}

\begin{document}

\begin{titlepage}
\thispagestyle{empty}
\begin{center}
МИНОБРНАУКИ РОССИИ\\
Федеральное государственное бюджетное образовательное учреждение высшего образования\\
«Ярославский государственный университет им. П.Г. Демидова»\\
Кафедра математического моделирования
\end{center}

\vspace{0.8cm}
\begin{flushright}
Сдано на кафедру\\
«\underline{\hspace{1cm}}» \underline{\hspace{3cm}} 2026 г.\\
Заведующий кафедрой\\
к. ф.-м. н., доцент\\
\underline{\hspace{4cm}} И. С. Кащенко
\end{flushright}

\vspace{1.2cm}
\begin{center}
{Выпускная квалификационная работа}\\[0.5cm]
{\bfseries Разработка концепции физического генератора случайных чисел}\\[0.5cm]
специальность 01.03.02 Прикладная математика и информатика
\end{center}

\vspace{1cm}
\begin{flushright}
Научный руководитель\\
старший преподаватель\\
\underline{\hspace{3cm}} / Д. А. Савинов\\
«\underline{\hspace{1cm}}» \underline{\hspace{3cm}} 2026 г.\\[0.7cm]
Студент группы ПМИ-41БО\\
\underline{\hspace{3cm}} / В. Д. Шестаков\\
«\underline{\hspace{1cm}}» \underline{\hspace{3cm}} 2026 г.
\end{flushright}

\vfill
\begin{center}
Ярославль 2026 г.
\end{center}
\end{titlepage}

\tableofcontents
\newpage
\end{document}